% Generated draft: requires mathematical and source review.
\documentclass[11pt]{article}
\usepackage[margin=25mm]{geometry}
\usepackage{fontspec}
\setmainfont{Latin Modern Roman}
\usepackage{amsmath,amssymb,mathtools,unicode-math}
\setmathfont{Latin Modern Math}
\usepackage{enumitem,needspace,xurl,hyperref}
\setlength{\parindent}{0pt}
\setlength{\parskip}{5pt}
\setlength{\emergencystretch}{3em}
\newcommand{\N}{\mathbb{N}}
\newcommand{\Z}{\mathbb{Z}}
\newcommand{\Q}{\mathbb{Q}}
\newcommand{\R}{\mathbb{R}}
\newcommand{\C}{\mathbb{C}}
\newcommand{\F}{\mathbb{F}}
\newcommand{\A}{\mathbb{A}}
\newcommand{\PP}{\mathbb{P}}
\newcommand{\E}{\mathbb{E}}
\newcommand{\eps}{\varepsilon}
\newcommand{\dd}{\,\mathrm{d}}
\begin{document}
\section*{UnsolvedMath ID: 30003568}
\textbf{Source problem code:} catalogue entry
\subsection{Definitions and mathematical statement}
Let \(X\) be a complex projective manifold of dimension \(n\). Under the usual convention allowing singularities, a codimension-one holomorphic foliation \(\mathcal F\) is a saturated coherent subsheaf \(T_{\mathcal F}\subseteq T_X\) of rank \(n-1\), closed under the Lie bracket. Saturation means that \(T_X/T_{\mathcal F}\) is torsion-free. On the locus where this subsheaf is a subbundle, its integral leaves have complex codimension one.

In the usual meromorphic interpretation, a transverse projective structure is represented by a \(\PP^1\)-bundle over \(X\), a flat meromorphic connection on that bundle, and a meromorphic section generically transverse to the connection's horizontal foliation, whose pullback induces \(\mathcal F\). Away from singularities and poles, this gives local first integrals with values in \(\PP^1\), related on overlaps by locally constant elements of \(\operatorname{PGL}_2(\C)\). The supplied record does not specify whether this meromorphic convention, including the allowance of poles, is intended.

For a projective surface \(Y\), a one-dimensional foliation \(\mathcal G\) has tangent sheaf of generic rank one. For a rational map \(\pi:X\dashrightarrow Y\), the equality \(\mathcal F=\pi^*\mathcal G\) means that, on a dense open subset where the map and both foliations are regular, the tangent distribution of \(\mathcal F\) is
\[
T_{\mathcal F}=(\dd\pi)^{-1}(T_{\mathcal G}),
\]
and that its saturated extension to \(X\) is the given foliation. This inverse image must have generic rank \(n-1\). The record does not explicitly require \(\pi\) to be dominant or \(Y\) to be smooth; these conditions must be checked against the intended source formulation.

\textbf{Supplied conjecture (Cerveau--Lins Neto).} Is it true that, for every complex projective manifold \(X\) and every codimension-one holomorphic foliation \(\mathcal F\) on \(X\), at least one of the following holds?
\begin{enumerate}
\item \(\mathcal F\) is transversely projective.
\item There exist a projective surface \(Y\), a one-dimensional foliation \(\mathcal G\) on \(Y\), and a rational map \(\pi:X\dashrightarrow Y\) such that
\[
\mathcal F=\pi^*\mathcal G.
\]
\end{enumerate}
The alternatives are not asserted to be mutually exclusive. The definitions above express a conventional reading; the singularity, meromorphic-structure, and pullback conventions remain subject to source verification and human mathematical review.
\subsection{Short English statement}
The supplied Cerveau--Lins Neto conjecture asks whether every codimension-one holomorphic foliation on a complex projective manifold either has a transverse projective structure or comes by rational pullback from a one-dimensional foliation on a projective surface. The record labels the conjecture open; this status and the precise geometric conventions require independent review.
\subsection{Sources}
\begin{itemize}
\item Supplied local UnsolvedMath registry record, permanent ID \(30003568\), titled \emph{Classification of Codimension-One Projective Foliations}; source problem code \texttt{catalogue entry}; collection: Oberwolfach Reports Open Problems. The record attributes the formulation to the Cerveau--Lins Neto conjecture, gives proposed year 2017, and supplies no individual proposer.
\item \emph{Komplexe Analysis (2018)}, Oberwolfach Reports, pp.\ 2427--2473, DOI \texttt{10.4171/owr/2017/39}: \url{https://doi.org/10.4171/owr/2017/39}. This is the supplied report citation and also a supplied literature source. The citation's 2018 label differs from the record's 2017 proposed year and the year in the DOI; these details are retained as supplied.
\item \emph{Foliations and the Cerveau Lins Neto conjecture}, supplied literature source: \url{https://www.ams.org/notices/202408/rnoti-p1025.pdf}.
\item Literature assessment dated exactly \texttt{2026-08-21}: the supplied curation labels the problem open and states that special cases and structural results are known while the stated dichotomy for all projective manifolds remains unresolved. This is supplied dated curation, not an independently verified assessment. The background expressly requests independent verification of time-sensitive status and cited claims.
\item The registry labels the statement \texttt{corrected\_verified} and reports that the official report was checked and the two alternatives and pullback data restored. These are supplied verification claims; no independent source consultation is asserted in this draft.
\end{itemize}
\end{document}
